% ===== PRÉAMBULE CANONIQUE — à coller VERBATIM en tête de CHAQUE .tex =====
\documentclass[11pt,a4paper]{article}
\usepackage[utf8]{inputenc}\usepackage[T1]{fontenc}\usepackage{lmodern}
\usepackage[french]{babel}
\usepackage{amsmath,amssymb,mathtools}\usepackage{icomma}
\usepackage[version=4]{mhchem}          % \ce{...} : équations & formules
\usepackage{siunitx}\sisetup{locale=FR} % \qty{}{}, \si{}, \ang{}
\usepackage{chemfig}                    % \chemfig{...} : structures développées
\usepackage{xcolor}\usepackage{colortbl}
\usepackage{tikz}\usepackage{pgfplots}\pgfplotsset{compat=1.18}
\usepackage[most]{tcolorbox}\usepackage{enumitem}\usepackage{array,booktabs}
\usepackage[a4paper,margin=2.2cm]{geometry}
\usetikzlibrary{arrows.meta,calc,angles,quotes,positioning,patterns,backgrounds,shapes.geometric}
\definecolor{primary}{RGB}{222,49,99}\definecolor{primarydark}{RGB}{150,22,55}
\definecolor{defcol}{RGB}{0,114,178}\definecolor{methcol}{RGB}{52,92,148}
\definecolor{excol}{RGB}{0,135,90}\definecolor{attncol}{RGB}{200,40,55}
\tcbset{boxbase/.style={enhanced,breakable,boxrule=0.9pt,arc=2.5pt,
  left=3mm,right=3mm,top=2.3mm,bottom=2.3mm,fonttitle=\bfseries,coltitle=white}}
% --- boîtes TUTORIEL (noms EXACTS, ne pas renommer) ---
\newtcolorbox{cle}[1][]{boxbase,colback=defcol!6,colframe=defcol,
  title={\textsc{À retenir}\ifx\relax#1\relax\relax\else\textmd{ : \itshape #1}\fi}}
\newtcolorbox{methode}[1][]{boxbase,colback=methcol!7,colframe=methcol,sharp corners,
  title={\textsc{Méthode}\ifx\relax#1\relax\relax\else\textmd{ --- \itshape #1}\fi}}
\newtcolorbox{exemplebox}[1][]{boxbase,colback=excol!5,colframe=excol,
  title={\textsc{Exemple}\ifx\relax#1\relax\relax\else\textmd{ : \itshape #1}\fi}}
\newtcolorbox{piege}[1][]{boxbase,colback=attncol!6,colframe=attncol,
  title={\textsc{Piège}\ifx\relax#1\relax\relax\else\textmd{ : \itshape #1}\fi}}
\newtcolorbox{remarque}[1][]{boxbase,colback=black!4,colframe=black!45,
  title={\textsc{Remarque}\ifx\relax#1\relax\relax\else\textmd{ : \itshape #1}\fi}}
% --- boîtes EXERCICE / CORRIGÉ (noms EXACTS, auto-comptées) ---
\newtcolorbox[auto counter]{exo}[1][]{boxbase,colback=primary!4,colframe=primary,
  title={\textsc{Exercice}~\thetcbcounter\ifx\relax#1\relax\relax\else\textmd{ --- \itshape #1}\fi}}
\newtcolorbox[auto counter]{corrige}[1][]{boxbase,colback=excol!4,colframe=excol,
  title={\textsc{Corrigé}~\thetcbcounter\ifx\relax#1\relax\relax\else\textmd{ --- \itshape #1}\fi}}
\newcommand{\R}{\mathbb{R}}\newcommand{\N}{\mathbb{N}}\newcommand{\Z}{\mathbb{Z}}\newcommand{\ds}{\displaystyle}
% (un \usepackage{fancyhdr}+en-tête est possible ; il est ignoré côté web)
% =========================================================================
\begin{document}
\sisetup{per-mode=symbol}
\begin{corrige}[doser les réactifs d'une synthèse]
Le coefficient de \ce{Al2(SO4)3} vaut $1$ ; on multiplie \qty{0.300}{\mole} par chaque coefficient.
\begin{enumerate}[label=\alph*)]
  \item $n(\ce{Al(OH)3})=0,300\times\dfrac{2}{1}=\qty{0.600}{\mole}$, d'où
        $m=0,600\times78=\qty{46.8}{\gram}.$
  \item $n(\ce{H2SO4})=0,300\times\dfrac{3}{1}=\qty{0.900}{\mole}$, d'où
        $m=0,900\times98=\qty{88.2}{\gram}.$
  \item $n(\ce{H2O})=0,300\times\dfrac{6}{1}=\qty{1.8}{\mole}.$
\end{enumerate}
\end{corrige}

\begin{corrige}[tableau d'avancement avec un réactif en excès]
\begin{enumerate}[label=\alph*)]
  \item $n(\ce{C2H5OH})=0,2-\xi$ ; \; $n(\ce{CO2})=2\xi$ ; \; $n(\ce{H2O})=3\xi$.
  \item L'oxygène étant en excès, c'est l'éthanol qui s'annule :
        $0,2-\xi=0\Rightarrow \xi_{\max}=\qty{0.2}{\mole}.$
  \item $n(\ce{CO2})=2\times0,2=\qty{0.4}{\mole}$, d'où $V=0,4\times22,4=\qty{8.96}{\litre}.$
\end{enumerate}
\end{corrige}

\begin{corrige}[combustion du kérosène : la chaîne complète]
\begin{enumerate}[label=\alph*)]
  \item $m=\rho V=0,763\times100=\qty{76.3}{\gram}$, puis
        $n(\ce{C14H30})=\dfrac{76,3}{198}\approx\qty{0.385}{\mole}.$
  \item $n(\ce{CO2})=n(\ce{C14H30})\times\dfrac{28}{2}=0,385\times14\approx\qty{5.39}{\mole}.$
  \item $V(\ce{CO2})=n\,V_{\mathrm{m}}=5,39\times22,4\approx\qty{121}{\litre}.$
\end{enumerate}
\end{corrige}

\begin{corrige}[dosage par titrage]
\begin{enumerate}[label=\alph*)]
  \item Les coefficients sont $1$ pour \ce{Cr2O7^2-} et $6$ pour \ce{Fe^2+}, donc à l'équivalence :
        $\dfrac{n(\ce{Cr2O7^2-})}{1}=\dfrac{n(\ce{Fe^2+})}{6}$, soit
        $n(\ce{Cr2O7^2-})=\dfrac{n(\ce{Fe^2+})}{6}.$
  \item $n(\ce{Fe^2+})=C\,V=0,60\times0,012=\qty{7.2e-3}{\mole}.$
  \item $n(\ce{Cr2O7^2-})=\dfrac{\num{7.2e-3}}{6}=\qty{1.2e-3}{\mole}$, d'où
        $[\ce{Cr2O7^2-}]=\dfrac{n}{V}=\dfrac{\num{1.2e-3}}{0,010}=\qty{0.12}{\mole\per\litre}.$
\end{enumerate}
\end{corrige}

\end{document}
